\chapter{Použitý framework aplikace}

\section{Operační systém}
V týmu je jako hlavní operační systém pro vývoj hry \emph{Surreal Bowling: Salvation Path} používán operační systém Linux. Bylo tak rozhodnuto, jelikož nabízí rychlejší workflow než Windows. Během práce také bylo zjištěno, že pokud na Unity pracují dva vývojáři, jeden na Linuxu a druhý na Windows, stává se, že se některé soubory v Unity projektu stanou poškozené. Vývoj videoher v operačním systému Linux je dnes proveditelný a programy jako Unity Engine, VS Code nebo Git fungují stejně dobře, ne-li lépe, než na Windows. \cite{linux-gamedev}

\section{Programování}

\subsection{C\#}
Pro psaní kódu byl použit multiplatformní programovací jazyk C\#, mimo jiné i nejpopulárnější jazyk pro .NET. \cite{csharp} Více o C\# v části \ref{sec:herni-engine}.

\subsection{Textový editor}
Co se textového editoru týče, byl použit VS Code z důvodu kompatibility s nejvíce pluginy třetích stran. \cite{vscode-vs-codium}

\subsection{Version Control}
Pro spolupráci ve vývojářském týmu byl použit Git – systém pro správu verzí, neboli „version control system“ (VCS). Přes Git jsou posílány lokální změny do online repozitáře projektu na webovou platformu GitHub, kde má přístup i náš mentor práce, aby měl přehled o tom, jak se s vývojem vede. \cite{git-intro} 

\subsection{AI (Umělá Inteligence)}
Pro usnadnění psaní kódu, provádění repetitivních úkolů nebo přesouvání větších objemů dat byly použity AI asistenty ChatGPT \cite{chatgpt} a Augment Code \cite{augment}. Později i lokální LLM (Large Language Model) v programu LM Studio \cite{lmstudio} (konkrétně qwen/qwen3.5-9b \cite{qwen}) integrovaný do nástroje OpenCode \cite{opencode}, což umožnilo daný model používat jako agenta v terminálu, který dokáže indexovat celé projekty a provádět v nich úpravy.

\section{Herní engine}
\label{sec:herni-engine}
Jakožto nejdůležitější software pro vývoj videohry dnes je herní engine. Pro tento projekt byl zvolen multiplatformní herní engine Unity vyvinutý společností Unity Technologies \cite{unity} z několika prostých důvodů. Jeden z nejhlavnějších je ten, že v Unity se defaultně programuje v programovacím jazyce C\#.

Dále byly také brány v potaz licence různých herních enginů. Mezi nejznámější patří právě Unity, Unreal Engine, Godot Engine nebo GameMaker. Unity prozatím stačilo jako nejrozumnější cesta, kde je možné videohry vydávat pod Unity Free License \cite{unity-pricing}, pokud jednotlivec nebo tým z výdělku z videohry nebo od sponzora přijímá méně než 200 000 amerických dolarů za posledních 12 měsíců. Jedna z největších nevýhod je, že při startu zadarmo vyvinuté videohry v Unity se při startu videohry objeví úvodní obrazovka (splash screen) s nápisem „Made with Unity“, což pro účely této práce zas tak nevadí.

Videohra vyvíjena v této práci také není nijak složitá v porovnání s dnešním videoherním standardem a Unity se obecně považuje jako příjemnější volba pro menší projekty \cite{unity-small-projects}, jako mohou být mobilní aplikace nebo právě jednoduché 2D projekty.

\section{Hudba a zvuk}

\subsection{FL Studio}
Pro tvorbu hudby a zvukových efektů byl použit program FL Studio.

FL Studio je software (Digital Audio Workstation) pro skládání, produkci a mix hudby za pomocí digitálních nástrojů, pluginů, samplů a MIDI signálů.

Mezi pluginy třetích stran, které byly frekventovaně používány pro tvorbu či modifikaci zvuku, patří LABS \cite{labs}, FLEX \cite{flex}, Sytrus \cite{sytrus} a Deelay \cite{deelay}. Plugin Deelay využívám především pro jeho pokročilé možnosti filtrace, časové a výškové modifikace a efekty jako Chamber nebo vrstvené komplexní delaye, které dodávají zvukům glitcheovou, zkorumpovanou atmosféru. \cite{flstudio}

\subsection{FMOD Studio}
\label{sec:fmod-studio}
Původně byly pro implementaci zvuku do této hry používány low-level nástroje přímo od Unity, například AudioSource, AudioListener a AudioMixer.

Později bylo implementováno FMOD Studio -- profesionální nástroj (audio middleware) pro úpravu a implementaci zvuku do videoher, vyvinutý společností Firelight Technologies \cite{fmod}. FMOD Studio umožňuje sound designerům vytvářet komplexní zvukové systémy pomocí přívětivého rozhraní bez nutnosti složitého programování, zatímco vývojáři mohou tyto zvuky snadno propojit s kódem pomocí Unity integrace. FMOD studio nabízí pokročilé funkce jako pocit 3D prostoru, automatizace zvuku nebo nejrůznější zvukové efekty (např. reverb).

\section{Ostatní}
Pro efektivní spolupráci v týmu je používán poznámkový software Obsidian \cite{obsidian}, kde předplatné Obsidian Sync zefektivňuje práci jak každodenně, tak například i na game jamech. Toto předplatné umožňuje ve stejný čas přistupovat ke všem informacím, nápadům, řešením, skriptům a zapisovat nové, aniž by bylo třeba cokoli manuálně přeposílat přes komunikační médium.

Pro komunikaci je v týmu používána platforma Discord \cite{discord} v založeném serveru pod názvem \emph{Still Thinking Studio}, na kterém jsou všichni členové týmu i spolupracující zvenčí. Discord umožňuje psát si jako v každé běžné chatovací aplikaci, spravovat desítky různých textových a hovorových kanálů (oprávnění, přístup, kategorie…) a rychle si posílat obrázky, zvukové soubory nebo videa pro nápady a zpětnou vazbu od ostatních členů týmu.